% Part: many-valued-logic
% Chapter: three-valued-logics
% Section: kleene

\documentclass[../../../include/open-logic-section]{subfiles}

\begin{document}

\olfileid{mvl}{thr}{skl}

\olsection{ક્લીની તર્કશાસ્ત્રો}

સ્ટીફન ક્લીનીએ બે ત્રિમૂલ્ય તર્કશાસ્ત્રો રજૂ કર્યાં. તેમની પ્રેરણા
એવા તર્કશાસ્ત્રમાંથી આવી જેમાં સત્યમૂલ્યોને સંગણનાત્મક પ્રક્રિયાઓના
પરિણામો માનવામાં આવે છે: કોઈ પ્રક્રિયા $\True$ અથવા $\False$ આપી
શકે, પણ તે અટકે જ નહિ એવું પણ બને. ત્યારે તેને અનુરૂપ સત્યમૂલ્ય
અવ્યાખ્યાયિત હોય છે અને તેને~$\Undef$ વડે દર્શાવીએ છીએ.

વિધાન~$!A$નો નિષેધ ગણવા માટે પહેલાં~$!A$નું મૂલ્ય ગણવું પડે અને
પછી તેનું વિપરીત મૂલ્ય આપવું પડે. જો $!A$ની ગણતરી અટકે નહિ, તો
આખી પ્રક્રિયા પણ અટકતી નથી; તેથી $\Undef$નો નિષેધ $\Undef$ છે.

સંયોજન $!A \land !B$ ગણવાની બે રીતો છે. પહેલાં~$!A$ અને પછી
$!B$ ગણીએ, તો બંનેનાં પરિણામ~$\True$ હોય ત્યારે જવાબ $\True$,
અન્યથા $\False$ મળે. બેમાંથી એક પણ ગણતરી અટકે નહિ, તો આખી
પ્રક્રિયા પણ અટકતી નથી. તેથી આ પદ્ધતિમાં એક ઘટક અવ્યાખ્યાયિત
હોય, તો સંયોજન પણ અવ્યાખ્યાયિત છે. વિયોજન માટે પણ એવું જ છે.

પરંતુ $!A$ અને $!B$ને સમાંતરે ગણાવી શકીએ તો વધુ સારું થાય.
બેમાંથી એક પ્રક્રિયા અટકી $\False$ આપે, તો આપણે અટકી શકીએ,
કારણ કે જવાબ ખોટો જ હશે. એટલે આ કિસ્સામાં એક ઘટક ખોટો હોય
તો બીજો અવ્યાખ્યાયિત હોવા છતાં સંયોજન ખોટું છે. એવી જ રીતે,
વિયોજન સમાંતરે ગણતાં કોઈ એક ઘટક સાચો નીકળે એટલે અટકી શકાય:
વિયોજન સાચું જ હશે. આમ એક ઘટક અવ્યાખ્યાયિત હોય તોપણ સંયુક્ત
વિધાનનું પરિણામ જાણી શકાય છે. આ અર્થઘટનમાં $\Undef$ને
``અવ્યાખ્યાયિત'' કરતાં ``અજ્ઞાત'' તરીકે વાંચી શકીએ.

આ બે અર્થઘટનોમાંથી ક્લીનીનાં પ્રબળ અને દુર્બળ તર્કશાસ્ત્રો મળે છે.
પ્રેરણ સંયોજકને $\lnot !A \lor !B$ની સમકક્ષ તરીકે વ્યાખ્યાયિત કરીએ છીએ.

\begin{defn}
\emph{પ્રબળ ક્લીની તર્કશાસ્ત્ર}~$\LogKs$ નીચેની શ્રેણિક વડે વ્યાખ્યાયિત થાય છે:
\begin{enumerate}
  \item $\lnot$, $\land$, $\lor$, $\lif$ ધરાવતી અને
  અસત્ય-અચલ વિનાની પ્રમાણભૂત વિધાનાત્મક ભાષા~$\Lang L_0$ની ઉપભાષા.
  \item સત્યમૂલ્યોનો ગણ $V = \{\True, \Undef, \False\}$.
  \item એકમાત્ર નિર્દિષ્ટ મૂલ્ય $\True$ છે; એટલે $V^+ = \{\True\}$.
  \item સત્યવિધેયો નીચેની સારણીઓ વડે અપાય છે:
  \begin{center}
    \begin{tabular}{c|c}
      $\tf{\lnot}$ & \\
      \hline
      $\True$ & $\False$ \\
      $\Undef$ & $\Undef$ \\
      $\False$ & $\True$
    \end{tabular}
    \quad
    \begin{tabular}{c|ccc}
      $\tf{\land}[\LogKs]$ & $\True$ & $\Undef$ & $\False$ \\
      \hline
      $\True$ & $\True$ & $\Undef$ & $\False$ \\
      $\Undef$ & $\Undef$ & $\Undef$ & $\False$\\
      $\False$ & $\False$ & $\False$ & $\False$
    \end{tabular}
    \\[2ex]
    \begin{tabular}{c|ccc}
      $\tf{\lor}[\LogKs]$ & $\True$ & $\Undef$ & $\False$ \\
      \hline
      $\True$ & $\True$ & $\True$ & $\True$ \\
      $\Undef$ & $\True$ & $\Undef$ & $\Undef$ \\
      $\False$ & $\True$ & $\Undef$ & $\False$
    \end{tabular}
    \quad
    \begin{tabular}{c|ccc}
      $\tf{\lif}[\LogKs]$ & $\True$ & $\Undef$ & $\False$ \\
      \hline
      $\True$ & $\True$ & $\Undef$ & $\False$ \\
      $\Undef$ & $\True$ & $\Undef$ & $\Undef$ \\
      $\False$ & $\True$ & $\True$ & $\True$
    \end{tabular}
  \end{center}
\end{enumerate}
\end{defn}

\begin{defn}
\emph{દુર્બળ ક્લીની તર્કશાસ્ત્ર}~$\LogKw$ નીચેની શ્રેણિક વડે વ્યાખ્યાયિત થાય છે:
\begin{enumerate}
  \item $\lnot$, $\land$, $\lor$, $\lif$ ધરાવતી અને
  અસત્ય-અચલ વિનાની પ્રમાણભૂત વિધાનાત્મક ભાષા~$\Lang L_0$ની ઉપભાષા.
  \item સત્યમૂલ્યોનો ગણ $V = \{\True, \Undef, \False\}$.
  \item એકમાત્ર નિર્દિષ્ટ મૂલ્ય $\True$ છે; એટલે $V^+ = \{\True\}$.
  \item સત્યવિધેયો નીચેની સારણીઓ વડે અપાય છે:
  \begin{center}
    \begin{tabular}{c|c}
      $\tf{\lnot}$ & \\
      \hline
      $\True$ & $\False$ \\
      $\Undef$ & $\Undef$ \\
      $\False$ & $\True$
    \end{tabular}
    \quad
    \begin{tabular}{c|ccc}
      $\tf{\land}[\LogKw]$ & $\True$ & $\Undef$ & $\False$ \\
      \hline
      $\True$ & $\True$ & $\Undef$ & $\False$ \\
      $\Undef$ & $\Undef$ & $\Undef$ & $\Undef$\\
      $\False$ & $\False$ & $\Undef$ & $\False$
    \end{tabular}
    \\[2ex]
    \begin{tabular}{c|ccc}
      $\tf{\lor}[\LogKw]$ & $\True$ & $\Undef$ & $\False$ \\
      \hline
      $\True$ & $\True$ & $\Undef$ & $\True$ \\
      $\Undef$ & $\Undef$ & $\Undef$ & $\Undef$ \\
      $\False$ & $\True$ & $\Undef$ & $\False$
    \end{tabular}
    \quad
    \begin{tabular}{c|ccc}
      $\tf{\lif}[\LogKw]$ & $\True$ & $\Undef$ & $\False$ \\
      \hline
      $\True$ & $\True$ & $\Undef$ & $\False$ \\
      $\Undef$ & $\Undef$ & $\Undef$ & $\Undef$ \\
      $\False$ & $\True$ & $\Undef$ & $\True$
    \end{tabular}
  \end{center}
\end{enumerate}
\end{defn}
\sourcecorrection{OLMV-007}{સ્થિર અંગ્રેજી મૂળ બંને
શ્રેણિકોમાં આખી પ્રમાણભૂત ભાષા~$\Lang L_0$ કહે છે;
પરંતુ અગાઉ તેની અંદર અસત્ય-અચલ પણ છે, જ્યારે અહીં
તેનું સત્યવિધેય અપાયું નથી. તેને અસત્ય મૂલ્ય આપીએ
તોપણ તેના નિષેધથી પુનરુક્તિ મળે છે અને નીચેનું
પ્રમેય ખોટું પડે છે. તેથી બંને શ્રેણિકોને અહીં
ચોક્કસ ચાર-સંયોજક, અચલ વિનાની ઉપભાષા સુધી
મર્યાદિત કરી છે.}

\begin{prop}
  $\LogKs$ અને $\LogKw$માં કોઈ પુનરુક્તિ નથી.
\end{prop}

\begin{proof}
  બધાં !!{propositional variable}s~$p$ માટે જો
  $\pAssign v(p) = \Undef$, તો કોઈ પણ સૂત્ર~$!A$નું સત્યમૂલ્ય
  $\pValue v(!A) = \Undef$ હશે, કારણ કે બંને તર્કશાસ્ત્રોમાં
  \[
    \tf{\lnot}(\Undef) = \tf{\lor}(\Undef, \Undef) = \tf{\land}(\Undef,
  \Undef) = \tf{\lif}(\Undef, \Undef) = \Undef
  \]
  અને $\LogKs$ કે $\LogKw$માંથી કોઈ પણ તર્કશાસ્ત્રમાં
  $\Undef \notin V^+$.
  તેથી આ !!{valuation} હેઠળ $!A$ નિર્દિષ્ટ નથી.
\end{proof}

પ્રબળ અને દુર્બળ ક્લીની તર્કશાસ્ત્રોમાં પુનરુક્તિ ન હોવા છતાં
તેમના ફલિતતા-સંબંધો ગૌણ નથી.

\begin{prob}
  નીચેનામાંથી કયા સંબંધો (a)~પ્રબળ અને (b)~દુર્બળ ક્લીની
  તર્કશાસ્ત્રમાં સાચા છે? દરેક માટે સત્યતાસારણી આપો.
  \begin{enumerate}
    \item $p, p \lif q \Entails q$
    \item $p \lor q, \lnot p \Entails q$
    \item $p \land q \Entails p$
    \item $p \Entails p \land p$
    \item $p \Entails p \lor q$
  \end{enumerate}
\end{prob}

દિમિત્રી બોચ્વારે $\Undef$નો અર્થ ``નિરર્થક'' લીધો અને
જૂઠ્ઠાણાના વિરોધાભાસ જેવા વિરોધાભાસો ઉકેલવા માટે વિરોધાભાસી
વિધાનોને~$\Undef$ મૂલ્ય આપવાની દરખાસ્ત કરી. તેણે એવો તર્કશાસ્ત્ર
રજૂ કર્યો જે મૂળભૂત રીતે વધારાના સંયોજકો સાથેનું દુર્બળ ક્લીની
તર્કશાસ્ત્ર છે. એમાંથી બે ``બાહ્ય નિષેધ'' અને ``અવ્યાખ્યાયિત છે''
કારક છે:
\begin{center}
  \begin{tabular}{c|c}
    $\tf{\sim}$ & \\
    \hline
    $\True$ & $\False$ \\
    $\Undef$ & $\True$ \\
    $\False$ & $\True$
  \end{tabular}
\quad
  \begin{tabular}{c|c}
    $\tf{+}$ & \\
    \hline
    $\True$ & $\False$ \\
    $\Undef$ & $\True$ \\
    $\False$ & $\False$
  \end{tabular}
\end{center}

\begin{prob}
  બોચ્વારના તર્કશાસ્ત્રમાં $\lnot$ અને~$+$ વડે $\sim$ વ્યાખ્યાયિત
  કરી શકાય છે? એટલે માત્ર !!{propositional
  variable}~$p$ ધરાવતું અને $\sim$ વિનાનું એવું !!a{formula}
  શોધો જેનું સત્યમૂલ્ય હંમેશાં~$\mathord{\sim}p$ જેટલું જ હોય.
  તમારો જવાબ સાચો છે તે દર્શાવવા સત્યતાસારણી આપો.
\end{prob}

\end{document}
